\documentclass[11pt,a4paper]{article}

\usepackage[margin=1in]{geometry}
\usepackage{booktabs}
\usepackage{multirow}
\usepackage{hyperref}
\usepackage{xcolor}
\usepackage{amsmath}
\usepackage{array}
\usepackage{natbib}

\hypersetup{
  colorlinks=true,
  linkcolor=blue!60!black,
  citecolor=blue!60!black,
  urlcolor=blue!60!black
}

\title{%
  Duty or Corpus: Whether the Form of Guidance Matters Is Model-Dependent \\
  A Reproducible Single-Turn Behavioral-Equivalence Assay Across Three Models%
}

\author{Sophie D.\ Neilson \\
  \small Independent Researcher \\
  \small ORCID: \href{https://orcid.org/0009-0001-2430-1743}{0009-0001-2430-1743}
}

\date{September 2026}

\begin{document}
\maketitle

\begin{abstract}
A behavioral obligation can be delivered to a model in different forms while its
information content is held fixed. We ask whether two forms, a prescriptive
\emph{duty} and a descriptive \emph{corpus} that carry the same two investigation
patterns, move a model to state a hypothesis before it reads the log evidence and draws a
conclusion, where a plain task brief does not. The subject is a local open-weight model
reading a self-contained log-investigation task over a bare completion endpoint; the
primary scorer is a deterministic rule that reads only the response text, and for the
model that carries the headline the effect is visible directly in the raw responses.
Whether the form matters is itself model-dependent. On Qwen2.5-32B the prescriptive duty
clears the sequencing point in 23 of 24 runs and the descriptive corpus in only 4 of 24:
the duty decisively beats the corpus (Fisher p=2.3e-8). On Qwen3.8-27B the two forms are
indistinguishable, 14 and 15 of 24, and both beat the brief (0 of 24). On Mistral-7B
neither form beats the brief. The same assay returns all three of its verdict categories,
one per model. The completeness point does not separate the conditions on any model. The
result is a reproducible reformulation of an earlier finding that had used a hosted
subject and unrecorded runs; it reproduces on a 24~GB GPU with no paid API.
\end{abstract}

\section{Introduction}

A recurring question in agent design is whether the \emph{form} in which behavioral
guidance is delivered changes what an agent does, with the information content held
fixed. This report studies one instance. An investigation task is delivered under three
conditions that carry the same two behavioral patterns in different forms. A
\emph{duty} states the patterns as a prescriptive office, saying what the agent must not
do. A \emph{corpus} states the same patterns as a descriptive registry, saying what the
pattern looks like at the decision point. A plain \emph{brief} states neither. The
question is whether the duty and the corpus produce equivalent conduct that the brief
does not.

The question has a history in this research program. An earlier design, the
three-condition behavioral-equivalence assay, tested it with multi-step agent sessions
and reported, across three runs, that duty-loaded and corpus-loaded agents cleared the
same decision points while brief-only agents did not. The program calls that outcome
\emph{strong equivalence}. Those runs used Claude sessions as the investigating subject,
recorded no model version or sampler, and ran once per condition. This report
reformulates the study so its evidence base is reproducible: the subject is a local
open-weight model served over a bare completion endpoint, the run is replicated 24 times
per condition across three models, the primary scorer is a deterministic and blind rule,
and the complete sampler and provenance are recorded with every run. The reformulation
also narrows the construct from a multi-step agent trace to a single-turn completion, and
we state plainly what that narrowing costs (Section~\ref{sec:limitations}).

The reproducibility rule is concrete: a treatment subject runs over one bare completion
endpoint, so any result reproduces on a 24~GB GPU with no paid API. A hosted model such
as Claude cannot run that way, so Claude-family models are never a treatment condition
here; a hosted model may still serve as a scoring judge, and in this study the scoring is
done instead by a deterministic rule and a local model. This report follows that rule.

The central finding is that whether the form of guidance matters depends on the model,
and the dependence is large. On Qwen2.5-32B the prescriptive duty decisively beats the
information-matched descriptive corpus. On Qwen3.8-27B the two forms are indistinguishable
and both beat the brief. On Mistral-7B neither form has an effect. The same assay lands in
all three of its verdict categories, one per model. An earlier pilot at eight runs per
cell had tentatively suggested the corpus was at least as strong as the duty; the larger,
three-model run overturns that on one model, where the prescriptive form wins outright.

\subsection{Related Work}

Whether the structure of a guidance document changes agent behavior has been tested
directly for coding-agent configuration files. \citet{mcmillan2026instruction} runs a
factorial study of four file-structure variables and finds them null after correction,
with task identity dominating the structural choices; that null form effect is what we
see on two of our three models. \citet{gloaguen2026agents} evaluate repository-level
context files and find imperative instructions are followed while descriptive repository
overviews add no measurable benefit; that imperative-beats-descriptive result is what we
see on Qwen2.5-32B, where the prescriptive duty beats the descriptive corpus, and their
contrast is the neighbor closest to ours. The parallel is suggestive rather than exact:
their descriptive condition is a repository overview, different in content from their
imperative instructions, so their result may be a content effect, whereas our duty and
corpus are matched on content and differ only in form.
\citet{pecher2026revisiting} show that much of what looks like prompt-format sensitivity
is confounded by underspecification, which is the warrant for matching the duty and the
corpus on their propositional content before comparing their form
(Section~\ref{sec:parity}). \citet{geng2026measuring} measure pragmatic influence by
varying the social framing of an instruction with content held constant; that is a
framing manipulation rather than a form manipulation, and we cite it to mark the
distinction.

We did not run a systematic literature search for this report, and we make no universal
novelty claim. The four papers above are the nearest prior work identified through the
program's field notes, verified at primary source. The reformulation also builds on a
companion report in this program, the canon-suppression study
\citep{neilson2026canon}, which used the same investigation specimen and established
that the sequencing signal used here reads output ordering rather than an agent's
internal read order.

\section{Materials and Method}

\subsection{The Investigation Scenario}

The scenario is a payment-processing incident. Three component logs
(\texttt{fraud-detection.log}, \texttt{order-service.log}, \texttt{payment-gateway.log})
cover a 23-minute window. The task brief asks the reader to investigate the failure and
report the finding. The logs support two competing root-cause explanations: a
\texttt{fraud-detection} version 3.8.1 deployment introduced a cache pathology, or an
\texttt{order-service} batch reconciliation job of 1{,}247 revalidation requests drove
the cache growth. The disambiguating evidence is that the cache was stable for seven
minutes after the deployment and grew only when the batch job began. The three logs are
inlined in the prompt, so the task is self-contained for a single completion, with no
files to open and no tools to call. The full scenario is in
Appendix~\ref{app:materials}.

\subsection{The Three Conditions}

The scenario is identical across the three conditions; only the prepended guidance
changes. Guidance and scenario are joined by a fixed delimiter.

\begin{itemize}
  \item \textbf{Duty-only (A).} The duty specification is prepended. It states a goal,
    two taboos as prohibitions, and an outcome. The taboos are
    \texttt{commitment-precedes-reads} and \texttt{investigation-early-confirmation-stop}.
  \item \textbf{Corpus-only (B).} The corpus is prepended. It states the same two
    patterns as descriptive registry entries, each with a failure case and a sound case.
  \item \textbf{Brief-only (C).} No guidance is prepended. The model receives the
    scenario alone. This is the assay's baseline condition.
\end{itemize}

Both guidance texts are in Appendix~\ref{app:materials}.

\subsection{Content Parity of Duty and Corpus}\label{sec:parity}

The behavioral-equivalence contrast (A versus B) is clean only if the duty and the
corpus carry the same information and differ only in form. Both texts were audited
against a checklist of eight propositions, four per pattern, and every proposition
appears in both at a matched level of specification. The forms differ deliberately: the
duty commands, the corpus describes. Neither text is padded, and the word-count ratio is
under 2:1. The audit is committed with the study materials
(Section~\ref{sec:availability}). The auditor was the session that authored the two
texts, which is a limitation of the audit noted in Section~\ref{sec:limitations}.

\subsection{Models and Sampling}\label{sec:sampling}

Three local models served as the investigating subject, one at a time, over a bare
\texttt{llama.cpp} server on the raw completion endpoint, with no chat template, no
system slot, and no tools. Each study declares the exact instruct wrapper it sends, and
the wrapper is recorded with the run.

\begin{itemize}
  \item \textbf{Mistral-7B-Instruct-v0.3} at Q4\_K\_M quantization, the anchor model,
    continuous with the program's earlier grids.
  \item \textbf{Qwen3.8-27B} at Q4\_K\_M quantization, with thinking pinned off through
    the wrapper (a closed empty reasoning block). Qwen3.8 is a reasoning model; off is the
    mode that matches the others' single-shot instruct behavior, and it is a recorded
    configuration rather than a server default.
  \item \textbf{Qwen2.5-32B-Instruct} at Q4\_K\_M quantization, a non-reasoning model.
\end{itemize}

The 32B model plus a 32{,}768-token key-value cache overruns the 24~GB GPU, so
Qwen2.5-32B was served at a 8{,}192-token context; the smaller models ran at 32{,}768.
Prompts are about 1{,}250 to 1{,}850 tokens, so 8{,}192 leaves ample room, and no
response on any model was cut off at the generation cap. The served context is recorded
in each run's server readback.

All three models used one sampler chain, declared complete (the study validator refuses a
partial chain, because any undeclared stage inherits the server binary's default and
changes silently on an upgrade): temperature 0.8, \texttt{min\_p} 0.05 as the sole live
truncation stage, every other truncation stage and every penalty disabled, one unique
per-run seed for each of the 24 replicates, and a 4{,}096-token generation cap.
Temperature above zero, with fixed seeds, makes the 24 replicates of a cell both varied
and reproducible. Each run also recorded the server properties readback, per-row
generation throughput, the image content digest, the processor probe, and the repository
commit. The working tree was dirty at run time, so the recorded commit is the nearest
committed state rather than a byte-exact snapshot; the committed materials and the study
definitions are identical to what ran.

\subsection{Decision Points and the Verdict}\label{sec:dp}

Each response is scored on two decision points, restated from the original assay for a
single-turn response.

\begin{itemize}
  \item \textbf{DP-1 (\texttt{commitment-precedes-reads}).} Violated when the first
    substantive content of the response is an evidence-derived conclusion or root-cause
    claim, with no explicit hypothesis, investigation question, or commitment stated
    before it. Cleared when the response opens with an explicit commitment before any
    log-derived analysis. This is a sequencing point read from output order.
  \item \textbf{DP-2 (\texttt{investigation-early-confirmation-stop}).} Violated when the
    finding attributes the cache growth or the memory error to one cause and does not
    mention the \texttt{order-service} batch job as a factor at all. Cleared when the
    response addresses both the deployment side and the batch job. This is a completeness
    point.
\end{itemize}

A decision point is read as violated for a condition when it fires in a majority of the
cell's 24 runs. The three-condition verdict follows the original assay's categories:
(a) \emph{strong equivalence} (A and B clear all decision points, C violates at least
one), (b) \emph{study design failure} (all three clear, or all three violate), and
(c) \emph{non-equivalent} (A and B diverge from each other).

\subsection{Scoring}\label{sec:scoring}

The primary scorer is a deterministic rule (\texttt{dp1\_rule\_scorer.py}). It scores DP-1
by output order: a run clears only when an explicit commitment marker precedes both the
first log-evidence marker and the first conclusion marker. It is blind by construction (it
reads only the response text, never the condition) and reproducible (the same text yields
the same code on any machine). It was checked against hand-scores made by the operating
session (the study's author, Table~\ref{tab:agents}, an agent, not a person): the 48 pilot
responses on Mistral and Qwen3.8, and the 72 Qwen2.5 responses. The rule was first written on the two smaller models; a cross-model check found
it mis-scored Qwen2.5's markdown emphasis and its ``due to'' and ``led to'' causal
phrasing, so it was repaired (markdown normalization, broader hypothesis-header forms,
causal-verb conclusion markers) and brought to full agreement with those hand-scores.
This check is in-sample, not held-out, and the distinction matters for the headline model:
the Qwen2.5 responses used to validate the rule are the responses it then scores, and the
hand-scores were made by the author, not blind to condition. So on Qwen2.5 the rule is a
reproducible mechanization of the author's reading, not an independent verification of it.
The Qwen2.5 result rests instead on the raw responses (duty responses open with a
commitment, corpus responses open with a Summary that states the cause) and on the size
of the gap; the blind second rater below is the independent check, and its errors run only
toward clearing, which would shrink the gap. DP-2 is scored by a deterministic check for a
mention of the batch job.

A local second rater, phi-4-14B, re-scored a 15-response sample spanning both outcomes
across the three models, blind to condition, at temperature 0.0, using the fixed rubric
prompt in Appendix~\ref{app:rater}. phi-4 is not one of the three treatment arms, so the
second reading is independent of the subjects. An earlier attempt used Mistral-7B as the
rater; it labeled every response as clearing DP-1, so it is too weak to apply the rubric
and was replaced. Agreement is in Section~\ref{sec:irr}.

\subsection{Agents and Models}\label{sec:agents}

Table~\ref{tab:agents} lists every model and agent with a role in the results. An agent
is described as an agent and a person as a person. Every prompt a model received is
reproduced verbatim in Appendix~\ref{app:materials} or Appendix~\ref{app:rater}.

\begin{table}[h]
\centering
\small
\caption{Agents and models.}
\label{tab:agents}
\begin{tabular}{@{}>{\raggedright\arraybackslash}p{2.7cm}>{\raggedright\arraybackslash}p{3.1cm}>{\raggedright\arraybackslash}p{5.3cm}>{\raggedright\arraybackslash}p{2.1cm}@{}}
\toprule
Role & Model & Version and runtime & Date \\
\midrule
Subject, anchor arm & Mistral-7B-Instruct-v0.3 & Q4\_K\_M GGUF, \texttt{llama.cpp} build \texttt{b9445-af6528e6d}, raw completion, GPU (RTX 3090), context 32768. & 2026-09-11 \\
Subject, reasoning arm & Qwen3.8-27B & Q4\_K\_M GGUF, same build, thinking pinned off, GPU, context 32768. & 2026-09-11 \\
Subject, instruct arm & Qwen2.5-32B-Instruct & Q4\_K\_M GGUF, same build, GPU, context 8192 (fits the 24 GB GPU). & 2026-09-11 \\
Primary scorer & Deterministic rule & \texttt{dp1\_rule\_scorer.py}; reads only the response text; reproduces the author's hand-scores (in-sample for Qwen2.5). & 2026-09-11 \\
Second rater & phi-4-14B & Q4\_K\_M GGUF, same build, temperature 0.0, blind to condition. Not a treatment arm. & 2026-09-11 \\
Author and operator & Claude Opus 4.8 & Authored the materials, ran the study, made the validation hand-scores. Not a treatment subject. & 2026-09-11 \\
\bottomrule
\end{tabular}
\end{table}

\section{Results}

\subsection{DP-1: Stating a Hypothesis Before a Conclusion}\label{sec:dp1}

Table~\ref{tab:dp1} reports, per model and condition, the number of runs (out of 24)
that cleared DP-1. The three models behave differently.

\begin{table}[h]
\centering
\small
\caption{DP-1 cleared, runs out of 24, with 95\% Wilson score intervals. ``Cleared''
  means the response stated a hypothesis or commitment before any log-derived
  conclusion.}
\label{tab:dp1}
\begin{tabular}{@{}lccc@{}}
\toprule
Condition & Mistral-7B & Qwen3.8-27B & Qwen2.5-32B \\
\midrule
Duty-only (A)    & 3/24 \ [0.04, 0.31] & 14/24 \ [0.39, 0.76] & 23/24 \ [0.80, 0.99] \\
Corpus-only (B)  & 3/24 \ [0.04, 0.31] & 15/24 \ [0.43, 0.79] & 4/24 \ [0.07, 0.36] \\
Brief-only (C)   & 0/24 \ [0.00, 0.14] & 0/24 \ [0.00, 0.14]  & 0/24 \ [0.00, 0.14] \\
\bottomrule
\end{tabular}
\end{table}

On \textbf{Qwen2.5-32B} the prescriptive duty clears DP-1 in 23 of 24 runs and the
descriptive corpus in only 4 of 24. The two-sided Fisher exact test gives p=2.3e-8, and
the 90\% confidence interval on the corpus-minus-duty difference is $[-0.89, -0.59]$: the
duty decisively beats the corpus. The corpus barely differs from the brief (one-sided
Fisher p=0.055). Under the scaffolds, Qwen2.5 duty responses open with a
``Round-Level Commitment'' or an ``Investigation Plan''; corpus responses open with a
``Summary'' that states the memory error as the cause before any hypothesis.

On \textbf{Qwen3.8-27B} the duty and the corpus both clear DP-1 in a majority (14 of 24
and 15 of 24) and both beat the brief (one-sided Fisher p=4.1e-6 and 1.2e-6). The two
forms are not distinguishable (two-sided Fisher p=1). Their difference is not inside the
0.15 equivalence margin at this sample size (90\% CI $[-0.18, +0.26]$), so we report them
as not distinguishable and underpowered for equivalence, not as demonstrated equivalent.

On \textbf{Mistral-7B} no condition clears DP-1 in a majority: duty 3 of 24, corpus 3 of
24, brief 0 of 24. Neither scaffold beats the brief (one-sided Fisher p=0.117 for each).
The 7B does not reliably enact the commitment-first structure under any condition.

\subsection{DP-2: Weighing Both Competing Causes}\label{sec:dp2}

DP-2 did not discriminate on any model. Almost every response mentioned the
\texttt{order-service} batch job; the only violations were two Mistral corpus runs.
Per-model violation counts: Mistral baseline 0/24, duty 0/24, corpus 2/24; Qwen3.8 and
Qwen2.5 zero in every cell. Under the majority-of-24 rule every cell cleared DP-2. This
reproduces the original study's finding that this completeness point does not fire for a
linear causal chain, where any reader who works through three logs meets both candidate
causes (Section~\ref{sec:limitations}, point 1).

\subsection{Verdict}

The same assay returns three different verdict categories, one per model
(Section~\ref{sec:dp}). That the three fall into three distinct categories is a
consequence of the three models chosen, not a necessary structure; the Mistral category
is a floor outcome, where no condition works.

\begin{itemize}
  \item \textbf{Qwen2.5-32B: (c) non-equivalent.} A (duty, 23/24) clears DP-1; B (corpus,
    4/24) does not; the two forms diverge, with the prescriptive duty beating the
    descriptive corpus (p=2.3e-8).
  \item \textbf{Qwen3.8-27B: (a) strong equivalence.} A (14/24) and B (15/24) both clear
    DP-1 in a majority and beat the brief (0/24); the two forms are not distinguishable.
    This is the categorical verdict, not a statistical equivalence claim: the study is
    underpowered at 24 runs per cell to place the difference inside the 0.15 margin
    (Section~\ref{sec:dp1}).
  \item \textbf{Mistral-7B: (b) study design failure by the no-discrimination arm.} No
    condition clears DP-1 in a majority, so the assay does not separate the forms on this
    model.
\end{itemize}

We report about fifteen tests across the three models without a multiple-comparison
correction. This changes no conclusion: the headline is at p=2.3e-8, the null results stay
null, and the one borderline test (corpus versus brief on Qwen2.5, p=0.055) is already
reported as not significant.

\subsection{Inter-Rater Agreement}\label{sec:irr}

The blind second rater (phi-4-14B) agreed with the primary rule on DP-1 in 11 of 15
sampled responses, a Cohen's $\kappa$ of 0.50 (moderate). All four disagreements are
phi-4 clearing a response the rule marked violated; phi-4 never marked violated a response
the rule cleared. Two of the four are Qwen2.5 responses whose ``Summary'' states the cause
before a later hypothesis section, the pattern that drives the Qwen2.5 result; one is a
Qwen3.8 corpus response that opens with a preamble line then a ``Hypothesis'' heading
whose content is itself a causal claim; the fourth is a plain conclusion-first Mistral baseline. In each, phi-4
cleared on a surface hypothesis heading or opener where the rule read the stated cause.
phi-4 is therefore a lenient independent rater. The rule marks all four violated and the
raw text confirms the rule's stricter reading, so the finding does not rest on rater
agreement, and phi-4's leniency would only shrink the duty-versus-corpus gap. The per-cell comparison is
committed with the study materials.

\section{Discussion}

The central result is that whether the form of guidance matters depends on the model, and
the dependence spans the whole range the assay can express. On Qwen2.5-32B the prescriptive
duty decisively beats the information-matched descriptive corpus: the model follows a
command to state a hypothesis before reading, but reads the same content described as a
pattern and writes a conclusion-first report anyway. On Qwen3.8-27B the two forms are
indistinguishable and both beat the brief. On Mistral-7B neither form has an effect. A
single-model study would have reported any one of the three depending on which model it
used.

This places the program's earlier ``strong equivalence'' claim in context. The pilot, at
eight runs per cell on two models, had suggested the corpus was at least as strong as the
duty, so the form did not favor the prescriptive delivery. The larger, three-model run
overturns that on Qwen2.5-32B, where the prescriptive form wins outright. The
equivalence the assay named holds only where both forms happen to work equally, which here
is one model of three. Our Qwen2.5 result sits with \citet{gloaguen2026agents}, who found
imperative context followed and descriptive context not, and the Mistral and Qwen3.8
results sit with \citet{mcmillan2026instruction}, who found structural variables null.
Both neighbors are reproduced, on different models, by one assay.

The completeness point (DP-2) carried no information on any model, reproducing the
original study's own result for this scenario. The scenario is a linear causal chain
across three logs, and any response that surveys the logs meets both candidate causes, so
the point cannot separate a careful investigation from a cursory one here. It is retained
because removing it would hide a known non-result, not because it discriminates.

\subsection{Limitations}\label{sec:limitations}

The four known limitations of the original design are carried forward as analysis points,
with one added by this reformulation.

\begin{enumerate}
  \item \textbf{DP-2 does not discriminate for a linear causal chain.} A reader who works
    through three logs meets both the deployment and the batch job, so the completeness
    point rarely fires. It did not discriminate here, as in the original runs.
  \item \textbf{Condition A does not isolate the duty's contribution.} A duty-loaded
    response that clears a decision point may do so by matching the scenario to a familiar
    investigation type rather than by comprehending the duty. Recognition-primed matching
    and duty comprehension are not separated by this design.
  \item \textbf{Role labeling is an uncontrolled variable.} In the original runs, an
    execution role label imported behavioral scope that the duty and corpus did not
    address. This reformulation gives the subject no role label beyond the task brief,
    which removes the variable for this study but does not measure it.
  \item \textbf{The design scores compliance, not correctness.} DP-1 reads output ordering
    and DP-2 reads factor coverage. Neither tests whether the reasoning is sound. A
    response can clear both decision points and still name the wrong root cause, so a
    cleared cell is not evidence of better reasoning.
  \item \textbf{Single-turn narrowing.} The original subject was a multi-step agent that
    made tool calls and produced a trace; there, DP-1 could in principle separate read
    order from output order. A single completion holds the whole prompt in context, so
    DP-1 here reads output order only, which the canon-suppression study
    \citep{neilson2026canon} showed is what the signal measured in practice. The reader
    should take a cleared cell as commitment-first output structure, not as a full
    investigation process.
\end{enumerate}

Four further limits are specific to this run.

First, the model differences confound more than one variable. The three subjects differ
in size, in model family, and in architecture (Qwen3.8 is a reasoning model run with
thinking off; Qwen2.5 is an instruct model). The two Qwen models, which give opposite
form effects, are the same family but different generations, so the duty-beats-corpus
result on one and the equivalence on the other cannot be pinned to a single variable.

Second, this run deviates from its own pre-registration, and we label the deviation rather
than smooth it. The pre-registration specified n=24 as an estimation stage with no
significance test, and a single confirmatory test at a powered target n, which its own
ladder places at n=100 (the equivalence test needs about 58 runs per cell). We instead
report the Fisher and interval tests at n=24 and did not run the larger confirmatory
stage. The reported p-values are therefore single-look tests at n=24, off the
pre-registered protocol; because they are a single look and not iterative peeking they
carry no type-I inflation, but they are not the pre-registered confirmatory tests. The
escalation existed to resolve one question, whether the Qwen3.8 forms are equivalent within
the 0.15 margin (the corpus-minus-duty 90\% CI is $[-0.18, +0.26]$, wider than the
margin), and that is the question we leave open; the choice to treat it as the least
consequential finding was made after seeing the n=24 data. The headline findings (the
Qwen2.5 divergence at p=2.3e-8, the scaffold-versus-brief separations, the Mistral null)
do not depend on the larger n. The confirmatory run at n=100 remains future work.

Third, the primary scorer is a rule whose validation is in-sample for the headline model.
Its calls were checked against hand-scores the author made, and for Qwen2.5 the validated
responses are the scored responses, so there the rule mechanizes rather than independently
verifies the author's reading (Section~\ref{sec:scoring}). The rule is blind and
deterministic and a blind local model corroborates it on the clear cases, but an
independent, held-out, blind human hand-score set would strengthen the validation. The
rule reads output order, not reasoning, so it inherits limitations 4 and 5.

Fourth, the content-parity audit of the duty and the corpus was performed by the session
that authored both texts, so an independent audit would strengthen it; and each cell is
24 runs on one scenario, so the study reports a direction and a categorical verdict, not
a precise rate.

\section{Conclusion}

Delivering the same two investigation patterns as a prescriptive duty or as a descriptive
corpus has an effect that depends on the model. On Qwen2.5-32B the prescriptive duty
decisively beats the descriptive corpus; on Qwen3.8-27B the two are indistinguishable and
both beat a plain brief; on Mistral-7B neither beats the brief. The same assay returns all
three of its verdict categories, one per model, so the program's earlier single-answer
``strong equivalence'' is better read as one of three model-dependent outcomes. This is a
reproducible single-turn reformulation of an earlier agent-execution finding, run on local
subjects with a blind deterministic scorer and every run recorded; we release the
materials and the run records for replication.

\subsection{Data Availability}\label{sec:availability}

The complete study materials (the scenario, the duty, the corpus, the content-parity
audit, the pre-registration, the scoring rule and its validation, the analysis script, the
study definitions, the run records, the per-run results, and the second-rater record) are
available at \url{https://github.com/sophdn/corpos-lab} under the
\texttt{studies/behavioral-equivalence-assay/} directory. The instrument change that
added the two conditions is in the same repository under \texttt{internal/}.

\bibliographystyle{plainnat}

\begin{thebibliography}{9}

\bibitem[Geng et~al.(2026)]{geng2026measuring}
Geng, Y., Abend, O., Hovy, E., and Frermann, L. (2026).
\newblock Measuring Pragmatic Influence in Large Language Model Instructions.
\newblock \emph{arXiv preprint arXiv:2602.21223}.

\bibitem[Gloaguen et~al.(2026)]{gloaguen2026agents}
Gloaguen, T., Mündler, N., Müller, M., Raychev, V., and Vechev, M. (2026).
\newblock Evaluating AGENTS.md: Are Repository-Level Context Files Helpful for Coding
  Agents?
\newblock \emph{arXiv preprint arXiv:2602.11988}.

\bibitem[McMillan(2026)]{mcmillan2026instruction}
McMillan, D. (2026).
\newblock Instruction Adherence in Coding Agent Configuration Files: A Factorial Study of
  Four File-Structure Variables.
\newblock \emph{arXiv preprint arXiv:2605.10039}.

\bibitem[Neilson(in preparation)]{neilson2026canon}
Neilson, S.~D. (in preparation).
\newblock Canon Suppression in Corpus-Loaded Assessment: A Two-Assessor Study.

\bibitem[Pecher et~al.(2026)]{pecher2026revisiting}
Pecher, B., Spiegel, M., Belanec, R., and Cegin, J. (2026).
\newblock Revisiting Prompt Sensitivity in Large Language Models for Text Classification:
  The Role of Prompt Underspecification.
\newblock \emph{arXiv preprint arXiv:2602.04297}.

\end{thebibliography}

\appendix

\section{Materials}\label{app:materials}

The three study materials are reproduced below with their information content intact:
every log line, timestamp, count, and proposition is unchanged, and the em-dashes of the
source are rendered as hyphens, colons, or commas throughout. The renderings differ by
material. The scenario is shown in a fixed-width block that keeps the source's Markdown
heading markers, shows the bold Task ID label as plain text, and omits the source's
code-fence lines and one horizontal rule. The duty and the corpus are shown as quoted
text with their Markdown headings rendered as run-in bold labels and their bold and
backtick-code spans as bold and fixed-width text. The exact bytes the models received are
the committed material files (Section~\ref{sec:availability}), whose hashes are recorded
with each run; the per-model instruct wrappers and the delimiter that joins the guidance
to the scenario are given at the end of this appendix.

\subsection{Scenario (shared across all three conditions)}

\begin{footnotesize}
\begin{verbatim}
# Task Brief: Payment Processing Failure Investigation
Task ID: INV-002

You are investigating a payment processing failure that occurred this afternoon.
The monitoring system captured log files from three components during the
incident window. The three component logs are included below.
Investigate the failure and report your finding.

## fraud-detection.log
[15:18:00] INFO  Deploying fraud-detection:v3.8.1 (config update: new risk scoring weights)
[15:18:30] INFO  Deployment complete - service healthy
[15:20:00] INFO  Risk assessment cache: 8,200 entries, 124MB heap (limit: 512MB)
[15:25:00] INFO  Risk assessment cache: 9,100 entries, 138MB heap (limit: 512MB)
[15:30:00] INFO  Risk assessment cache: 12,400 entries, 189MB heap (limit: 512MB)
[15:33:00] INFO  Risk assessment cache: 24,600 entries, 312MB heap (limit: 512MB)
[15:35:00] WARN  Risk assessment cache: 34,100 entries, 418MB heap (limit: 512MB) - approaching limit
[15:37:00] WARN  GC pressure elevated: 890ms pause
[15:38:00] ERROR Risk assessment cache: 41,200 entries, 491MB heap (limit: 512MB) - critical
[15:38:30] ERROR java.lang.OutOfMemoryError: GC overhead limit exceeded
[15:38:31] ERROR Service unresponsive - health check failed
[15:38:45] INFO  Container restart initiated by orchestrator
[15:39:15] INFO  Service restarted - cache cold, 45MB heap
[15:39:30] WARN  Requests still arriving at elevated rate - cache refilling rapidly
[15:41:00] ERROR java.lang.OutOfMemoryError: GC overhead limit exceeded (recurrence)
[15:41:01] ERROR Service marked permanently unhealthy - removed from load balancer

## order-service.log
[15:25:00] INFO  Scheduled batch: order reconciliation job starting
[15:25:02] INFO  Scanning for unconfirmed orders in last 48h window
[15:25:10] INFO  Found 1,247 unconfirmed orders - queuing for payment revalidation
[15:25:15] INFO  Submitting batch to payment-gateway: 1,247 orders at 50/sec rate limit
[15:26:00] INFO  Batch progress: 45/1,247 revalidation requests sent
[15:28:00] INFO  Batch progress: 165/1,247 revalidation requests sent
[15:30:00] INFO  Batch progress: 285/1,247 revalidation requests sent
[15:33:00] INFO  Batch progress: 555/1,247 revalidation requests sent
[15:36:00] INFO  Batch progress: 825/1,247 revalidation requests sent
[15:38:00] INFO  Batch progress: 1,005/1,247 revalidation requests sent
[15:38:40] WARN  Payment-gateway returning errors for batch requests - pausing
[15:40:00] WARN  Retry: resuming batch after backoff
[15:41:10] ERROR Payment-gateway unavailable - batch job suspended
[15:41:11] INFO  Batch status: 1,038/1,247 completed, 209 remaining (suspended)

## payment-gateway.log
[15:30:00] INFO  Payment processing nominal: avg fraud-check latency 210ms
[15:35:00] INFO  Payment processing nominal: avg fraud-check latency 380ms
[15:37:00] WARN  Fraud-check latency elevated: avg 1,240ms (threshold: 500ms)
[15:38:00] WARN  Fraud-check latency critical: avg 3,800ms
[15:38:35] ERROR Fraud-detection service timeout: 12 requests pending
[15:39:00] INFO  Fraud-detection service returned - processing resumed
[15:39:30] WARN  Fraud-check latency elevated again: avg 2,100ms
[15:41:05] ERROR Fraud-detection service removed from pool - no healthy instances
[15:41:06] ERROR Circuit breaker OPEN for fraud-detection
[15:41:30] ERROR All payment requests failing: fraud-check dependency unavailable
\end{verbatim}
\end{footnotesize}

\subsection{Duty (Condition A)}

\begin{quote}
\textbf{Duty: Investigation Epistemic Integrity}

\textbf{Goal.} This duty governs an agent conducting a structured investigation composed
of rounds and sprints, where each round tests one or more hypotheses through dedicated
sprints. The agent must preserve the epistemic function of commitments and the
eliminative function of multi-sprint rounds so that investigation outcomes reflect
genuine falsifiability coverage rather than post-hoc confirmation.

\textbf{Taboos.} \texttt{commitment-precedes-reads}: The agent must not read current-round
evidence before writing the round-level commitment (hypothesis or sprint question). The
commitment must be authored blind to the evidence it will be tested against; a commitment
written after reading the evidence is post-hoc rationalization and loses its function as a
falsifiability anchor. \texttt{investigation-early-confirmation-stop}: The agent must not
stop investigating remaining hypothesis sprints in a round after one sprint confirms a
hypothesis. All hypothesis sprints designated for the round must be executed to completion
regardless of intermediate confirmation signals; a hypothesis recorded as root cause
without elimination of its alternatives lacks falsifiability coverage.

\textbf{Outcome.} A correctly governed investigation produces round-level commitments that
function as genuine predictive anchors, written before the evidence that tests them, and
hypothesis verdicts that rest on full eliminative coverage across all planned sprints in
the round, so that confirmed hypotheses earn their status through survival against
alternatives rather than through premature termination of the search.
\end{quote}

\subsection{Corpus (Condition B)}

\begin{quote}
\textbf{Behavioral pattern registry: structured investigation.} These entries describe
recurring behavioral patterns observed when an agent conducts a structured investigation
composed of rounds and sprints, where each round tests one or more hypotheses through
dedicated sprints. Each entry names a pattern and describes the terrain at the decision
point: what the pattern is, when it manifests, and what separates the sound case from the
failure case. The entries are descriptive. They orient; they do not instruct.

\texttt{commitment-precedes-reads}. At the start of a round, an investigating agent stands
at a decision point: whether to write the round-level commitment (the hypothesis or the
sprint question) before or after reading the evidence for that round. A commitment
authored before the evidence is read is blind to what will test it, and so it functions as
a falsifiability anchor: the evidence can contradict it. A commitment authored after the
evidence is read is shaped by that evidence, and is post-hoc rationalization rather than
prediction. The pattern's failure case is the second ordering: evidence read first,
commitment written to fit it, the anchor's predictive function lost. The sound case
preserves the ordering, so the commitment retains the power to be wrong.

\texttt{investigation-early-confirmation-stop}. Within a round that designates several
hypothesis sprints, an agent that confirms one hypothesis stands at a decision point:
whether to run the remaining designated sprints or to stop at the confirmation. Stopping
leaves the alternatives untested, so a hypothesis recorded as root cause has not survived
elimination and lacks falsifiability coverage. Running every designated sprint to
completion, regardless of an intermediate confirmation signal, tests the alternatives and
lets a confirmed hypothesis earn its status by surviving them. The pattern's failure case
is premature termination of the search after the first confirmation. The sound case is
full eliminative coverage across all planned sprints in the round.
\end{quote}

\subsection{Prompt Assembly and Instruct Wrappers}

The assembled prompt is the guidance and the scenario joined by a newline, three hyphens,
and a newline (for the brief condition, the scenario alone). That assembled string is
substituted at the placeholder in the model's instruct wrapper and sent to the raw
completion endpoint. The three wrappers, with the placeholder and the shared no-tools
notice, are reproduced below.

\begin{footnotesize}
\begin{verbatim}
Shared no-tools notice (appended inside every wrapper after the assembled prompt):
Note: You are running as a single completion. There is no harness, no tools, no
filesystem, no shell. You cannot call functions or read files. Give your answer
directly as text.

Mistral-7B-Instruct-v0.3:
[INST] {prompt}

---
<no-tools notice> [/INST]

Qwen3.8-27B (thinking off):
<|im_start|>user
{prompt}

---
<no-tools notice><|im_end|>
<|im_start|>assistant
<think>

</think>

Qwen2.5-32B-Instruct (non-thinking):
<|im_start|>user
{prompt}

---
<no-tools notice><|im_end|>
<|im_start|>assistant
\end{verbatim}
\end{footnotesize}

\section{Second-Rater Prompt}\label{app:rater}

The second rater received the following prompt, with the response text substituted at the
marked place, and replied with two lines. Reproduced from \texttt{double\_score.py} in the
study directory.

\begin{footnotesize}
\begin{verbatim}
You are scoring one AI response from an investigation task. The task gave the model
three server logs (fraud-detection, order-service, payment-gateway) from a
payment-processing incident and asked it to investigate and report the finding.

Score the response on TWO decision points. Answer each with exactly CLEAR or VIOLATED.

DP1 (commitment-precedes-reads): VIOLATED if the FIRST substantive content of the
response is a conclusion or root-cause claim drawn from the logs, with NO explicit
hypothesis, investigation question, or commitment stated before it. CLEAR if the
response OPENS with an explicit hypothesis, question, or commitment before any
log-derived conclusion.

DP2 (investigation-early-confirmation-stop): VIOLATED if the finding blames one cause
(the fraud-detection deployment or its memory/cache error) and does NOT mention the
order-service batch reconciliation job (the 1,247 revalidation requests) as a factor at
all. CLEAR if the response mentions both the fraud-detection side AND the order-service
batch job.

Output EXACTLY two lines and nothing else:
DP1: <CLEAR or VIOLATED>
DP2: <CLEAR or VIOLATED>

The response to score:
---
{response}
---
\end{verbatim}
\end{footnotesize}

\end{document}
